\scanpage{31}
\mspage{10}
est ``transversal'' i.e. \deleted{\unsure{si}} $\exists$ voisinages cylindriques
de $X_{j'}$ \inserted{dans $X_j$} et un de $X_{k'}$ dans $X_k$, tels que l'application
$X_k\unsure{-|\to}X_j$ induit un \deleted{\unsure{hom.}}\inserted{morphisme} sur les
cylindres (compatible avec les structures
cylindriques) et de plus
$\overline E_{k',k}=\unsure{f_{jk}|p_{jk}}^{-1}(\overline E_{j',j})$.

Système de voisinages cylindriques \deleted{associé}
\inserted{(des $X_{j'}$ dans les $X_j$)}
compatible avec une stratification : tel que
dans b, les $E_{j'j}$ ($j$ fixé, $j'$ prédécesseur)
soient compatibles. Notion correspondante
pour les germes de voisinage cylindrique.

Stratification cylindrique : stratification munie
d'une telle ``structure cylindrique'', en
termes de germes de voisinages cylindriques.

\deleted{\uline{NB} \unsure{Une}}
\uline{NB} Une stratification précylindrique \inserted{resp. cylindrique}
est \unsure{sûrement}
``équisingulier'' au sens le plus fort qu'on
puisse introduire \unsure{qui} existe.

Mais j'ai oublié une condition, que j'exigerai
de tte stratification :
\[
(0)\qquad \forall i,j\in I,\quad\text{on a}\quad
X_i\cap X_j=\bigcup_{k\le i,j}X_k
\]
\margin{\uline{NB} si $i,j$ ont pour borne inf $h$,
cela implique $X_i\cap X_j=X_h$. \unsure{même pour} \illegible{3} \unsure{c)} $\unsure{X_i\cap X_j}$.}

Corollaire : Si $X_i\cap X_j\ne\varnothing$, alors $i,j$ ont un
minorant $h$ (tel que \deleted{\unsure{le}} $X_k\ne\varnothing$\ldots)
\margin{30}
